% Elaborado por Fabian Gomez
% Version 1.0

% =========================================================
% ICD de COMMS (Rover Olympus)
% =========================================================
\subsection{ICD de COMMS (Interfaz Rover–GCS y Relay)}
\label{sec:icd_comms}

\subsubsection{Alcance}
Este ICD define las interfaces \textbf{externas} y \textbf{lógicas} del subsistema COMMS del Rover Olympus
para el intercambio TM/CMD con la Estación Terrena (GCS) y, cuando aplique, a través de un \emph{nodo
retransmisor} (Relay) bajo un esquema \emph{store-and-forward}. Se cubren dos medios:
\begin{enumerate}
  \item \textbf{Wi-Fi 2.4\,GHz + UDP/IP + CSP} (canal de conveniencia/alta tasa en campo cercano).
  \item \textbf{UHF 430–440\,MHz + AX.25/KISS + CSP, GFSK 9600\,bps} (canal robusto de misión).
\end{enumerate}
\begin{sloppypar}
La arquitectura y requisitos vienen establecidos en el SyRS (§9.1.1/\allowbreak §9.4.1, RF‑006, COMM‑REQ‑001/\allowbreak 002)
y en la extensión COMMS UHF (COMM‑REQ‑003/\allowbreak 004/\allowbreak 005). El diseño práctico de UHF sigue la propuesta
del TFG ELANav Comms (KISS/\allowbreak AX.25, GFSK 9600, CSP/\allowbreak RDP y retransmisor).
\end{sloppypar}
% Referencias: SyRS (main.pdf) §9.1.1/§9.4.1 y TFG ELANav Comms

\subsubsection{Entidades y roles}
\begin{itemize}
  \item \textbf{Rover (C\&DH/COMMS):} origina/consume TM/CMD; expone dos interfaces de red: \emph{if\_udp} (Wi-Fi) e \emph{if\_kiss} (UHF).
  \item \textbf{GCS (Estación Terrena):} genera CMD y recibe TM; puede usar UDP/IP o radio UHF (vía TNC/Radio).
  \item \textbf{Relay (Retransmisor CSP):} recibe, almacena temporalmente y reenvía paquetes (store-and-forward) cuando no hay LoS directa.
\end{itemize}

\subsubsection{Interfaces físicas}
\paragraph{IF-COMMS-UDP (Rover$\leftrightarrow$GCS, Wi-Fi).}
\begin{itemize}
  \item Medio: 2.4\,GHz (IEEE\,802.11 b/g/n), canal/site \textbf{TBD}.
  \item Transporte: \textbf{UDP/IP}.
  \item Encapsulamiento de misión: \textbf{CSP}.
  \item Dirección IP Rover: \textbf{TBD} (DHCP o estática); Dirección IP GCS: \textbf{TBD}.
  \item Puertos UDP (CSP): \textbf{TBD} \emph{(RECOMENDADO:} rango 7000–10000 para evitar colisiones locales\emph{)}.
\end{itemize}
(\emph{Ref. SyRS §9.1.1/§9.4.1; COMM-REQ-001/002})%

\paragraph{IF-COMMS-KISS (Rover$\leftrightarrow$TNC, UHF).}
\begin{itemize}
  \item Enlace serie: USB CDC ACM $\rightarrow$ \emph{/dev/ttyACMx} (o UART nativa), \textbf{baudios TBD}
        \emph{(RECOMENDADO: 115200\,bps)}, 8N1.
  \item Protocolo: \textbf{KISS} para tramas \textbf{AX.25} (NRZI/CRC-16) moduladas \textbf{GFSK 9600\,bps}.
  \item TNC: \emph{NinoTNC 9600A} (modo 0000 = 9600\,GFSK AX.25). Nivel PTT/Audio según cableado del radio.
  \item Radio UHF (laboratorio): \emph{Kenwood TM-D710GA / TS-2000}; puerto de datos, PTT, entrada/salida audio.
  \item Antenas: \emph{Eggbeater} (omni, rover/relay) / \emph{Yagi} (direccional, estación). Coax y conectores \textbf{TBD}.
\end{itemize}
(\emph{Ref. TFG: KISS/AX.25, GFSK 9600, NinoTNC 9600A, Kenwood, Eggbeater/Yagi})%

\subsubsection{Interfaces lógicas y direccionamiento}
\paragraph{Pila de protocolos (ambos medios).}
\begin{center}
  \small \verb|Aplicación| $\rightarrow$ \textbf{CSP (RDP o no fiable)} $\rightarrow$
  \verb|UDP/IP| \emph{(Wi-Fi)} \textbf{o} \verb|KISS/AX.25| \emph{(UHF)} $\rightarrow$ \verb|RF|
\end{center}
(\emph{Ref. SyRS §9.4.1; TFG CSP sobre UDP o KISS/AX.25})%

\paragraph{IDs lógicos CSP.}
\begin{itemize}
  \item \textbf{Configurables} en despliegue. \emph{Ejemplo (TFG)}: \texttt{GCS=10}, \texttt{Rover=20}, \texttt{Relay=30}.
        Usar como \emph{valores de campaña}, documentados en el \emph{run card} del ensayo.
\end{itemize}
(\emph{Ref. TFG: ejemplo de direcciones CSP})%

\subsubsection{Nivel de implementación CSP y plan de migración}
\label{sssec:icd_comms_csp_level}

La implementación CSP del Rover Olympus se organiza en tres niveles de madurez.
La campaña del TFG opera en \textbf{Nivel~1}; los niveles superiores son trabajo futuro
cuando se integre el canal UHF.

\begin{sloppypar}
\begin{description}
  \item[Nivel~1 — empaquetado manual (implementado, campaña TFG)]
    El módulo \texttt{olympus\_hlc/\allowbreak csp.py} implementa el encabezado CSP de 32~bits
    mediante \texttt{struct.pack/\allowbreak unpack} de la biblioteca estándar de Python, sin
    dependencias externas. El campo CRC-32 se calcula con \texttt{zlib.crc32}
    (RF-006). La implementación cubre el canal Wi-Fi/UDP en campo cercano
    (COMM-REQ-001/002): envío de CMD (puerto~11), recepción de TM (puerto~10)
    y probe de reconexión \texttt{HB\_\allowbreak REQ} (puerto~1). \textbf{No soporta RDP ni
    routing multi-interfaz} — ambos innecesarios sobre WiFi local donde la
    pérdida de paquetes es despreciable.

  \item[Nivel~2 — libcsp con bindings Python (trabajo futuro, pre-UHF)]
    Migración de \texttt{csp.py} a \texttt{ctypes}/bindings oficiales sobre
    \texttt{libcsp.so}, compilado para RPi~5 con Yocto
    (\texttt{-{}-enable-python-\allowbreak bindings -\allowbreak-with-os=posix}).
    Habilita \textbf{RDP} (retransmisión con ACK) y tabla de rutas multi-interfaz.
    Requerido antes de integrar el canal UHF (COMM-REQ-003/004/005) porque RDP
    es la única capa que garantiza entrega sobre un enlace con pérdidas.
    El header CSP de 32~bits es idéntico al del Nivel~1;
    \texttt{GCSSource} y \texttt{gcs\_mock.py} no requieren cambios de lógica.

  \item[Nivel~3 — stack completo dual-homing (trabajo futuro, post-UHF)]
    Tabla de rutas libcsp con dos interfaces registradas:
    \texttt{CSP-UDP} (Wi-Fi, campo cercano) y
    \texttt{CSP-KISS} (UHF, operación robusta).
    La conmutación se produce automáticamente según disponibilidad del enlace
    y política RDP (\texttt{timeout\_\allowbreak RDP} / \texttt{max\_\allowbreak retries}).
\end{description}
\end{sloppypar}

\begin{sloppypar}
\textbf{Decisión de diseño:} permanecer en Nivel~1 durante el TFG minimiza
la deuda técnica de compilación cruzada en Yocto antes de completar la
validación del loop HLC$\leftrightarrow$LLC en hardware. La arquitectura es
extensible: agregar libcsp requiere portar únicamente \texttt{csp.py}
(90~líneas) sin tocar \texttt{GCSSource}, \texttt{HlcEngine} ni los monitores.
(\emph{Ref. COMM-REQ-001/002 cubiertos en Nivel~1; COMM-REQ-003/\allowbreak 004/\allowbreak 005 requieren Nivel~2+})
\end{sloppypar}

\subsubsection{Selección y conmutación de ruta}
La aplicación en C\&DH mantiene \textbf{dos interfaces CSP} activas: \emph{if\_udp} (Wi-Fi) y \emph{if\_kiss} (UHF).
La \textbf{política de ruta} se implementa en la tabla de enrutamiento CSP y la gestión de reintentos de \textbf{RDP}:
\begin{enumerate}
  \item \textbf{Ruta directa preferida}: Rover$\leftrightarrow$GCS (Wi-Fi o UHF directo).
  \item \textbf{Conmutación a Relay}: si se \emph{agotan reintentos/\allowbreak timeout RDP},
        se enruta Rover$\rightarrow$Relay$\rightarrow$GCS.
  \item \textbf{Retorno a directo}: al detectar \texttt{link\_restored}, se vuelve a la ruta directa.
\end{enumerate}
\begin{sloppypar}
Parámetros \textbf{TBD} (\emph{RECOMENDADO} para caracterización inicial, según TFG):
\emph{timeout\_\allowbreak RDP=500\,ms}, \emph{retries=3}.
(\emph{Ref.\ TFG: lógica y ejemplo de reintentos; SyRS §7.3.8 eventos
\texttt{link\_lost/\allowbreak link\_restored}}).
\end{sloppypar}

\subsubsection{Temporización y ventanas}
\paragraph{Wi-Fi (UDP/IP).}
Inyección de latencia \textbf{0–5\,s} (CON-002) \emph{solo} en el camino UDP/IP para emular retardo operacional.
(\emph{Ref. SyRS CON-002, §9.1.1})%

\paragraph{UHF (AX.25/KISS).}
Latencias y throughput dependen de ventana de aire y \textbf{RDP}. El TFG mostró que \textbf{RDP} garantiza
entrega (100\% en ensayos con 10/10 tramas) a costa de mayor tiempo total que el modo no fiable.%
(\emph{Ref. TFG, comparativa RDP vs no fiable})%

\subsubsection{Diccionario de mensajes (TM/CMD)}
\paragraph{Tipos TM/CMD (mínimos).}
\begin{itemize}
  \item \textbf{CMD de modo}: \verb|start/pause/homing/safe|; \textbf{CMD de configuración}: parámetros UC-01.
  \item \textbf{TM de estado (mínimo)}: \verb|mode|, \verb|timestamp|, \verb|link_status|, \verb|EPS.V|, \verb|EPS.I|.
\end{itemize}
La TM mínima se deriva de SyRS-030; extender con métricas UC-01 (slip, waypoint, etc.).%
(\emph{Ref. SyRS-030})%

\paragraph{Estructura lógica (CSP).}
\begin{itemize}
  \item \textbf{Encabezado CSP}: campos estándar (pri, src, dst, sport, dport, flags, etc.).
  \item \textbf{Payload}: TLV o binario compacto según tópico; \textbf{CRC/Checksum} según CSP.
\end{itemize}
(\emph{Ref. SyRS §9.4.1; uso de CSP como encapsulamiento})%

\paragraph{Encapsulamiento UHF.}
\begin{itemize}
  \item \textbf{KISS/AX.25}: trama AX.25 con \textbf{CRC-16}; KISS delimita por puerto serie (FEND/FESC).
  \item \textbf{GFSK 9600}: modulación en el TNC; PTT y audio según esquema del radio.
\end{itemize}
(\emph{Ref. TFG: AX.25/KISS, GFSK 9600})%

\subsubsection{Gestión de enlace y errores}
\paragraph{Detección de pérdida/recuperación.}
\begin{itemize}
  \item \texttt{link\_lost}: disparada por fallos de ACK RDP o heartbeats (Wi-Fi); transición a \emph{GestiónEnlace}.
  \item \texttt{link\_restored}: retorno a ruta directa y reanudación de TM periódica.
\end{itemize}
(\emph{Ref. SyRS §7.3.8 diccionario de eventos y estados})%

\paragraph{Política de reintentos (RDP).}
Parámetros \textbf{TBD} por campaña (timeout, max\_retries). \emph{RECOMENDADO} inicial: 500\,ms / 3 reintentos
como en el TFG de referencia para provocar conmutación a Relay en pérdida de LoS.%
(\emph{Ref. TFG})%

\subsubsection{Asignación de puertos/recursos (propuesta inicial)}
\paragraph{Wi-Fi / UDP.}
\begin{itemize}
  \item \texttt{GCS\_LISTEN\_PORT} (RPi\,5 escucha CMD): \textbf{9000}
        (\emph{resuelto en \texttt{config.py}}).
  \item \texttt{GCS\_REPLY\_PORT} (GCS escucha TM): \textbf{9001}
        (\emph{resuelto en \texttt{config.py}}).
  \item \texttt{IP\_ROVER / IP\_GCS}: \textbf{TBD} (DHCP o estática).
  \item Puertos CSP de servicio: CMD=11, TM=10, HB=1
        (\emph{resuelto en \texttt{config.py}: \texttt{CSP\_PORT\_CMD/TM/HB}}).
  \item Direcciones CSP de nodo: GCS=1, HLC=2
        (\emph{resuelto en \texttt{config.py}: \texttt{CSP\_ADDR\_GCS/HLC}}).
\end{itemize}

\paragraph{KISS / TNC.}
\begin{itemize}
  \item \texttt{/dev/ttyACMx} o \texttt{/dev/ttyUSBx}, \textbf{baudios TBD} (\emph{RECOMENDADO: 115200}).
  \item \texttt{KISS\_PORT\_ID}: \textbf{TBD} (si se multiplexan puertos KISS).
\end{itemize}

\subsubsection{Mapeo N2 y PBS/ASM}
Correspondencia con la matriz N2 y PBS/ASM:
\begin{itemize}
  \item \textbf{ASM-063} (TNC NinoTNC 9600A): [DATA] USB/UART (KISS), [CTRL] PTT, [PWR] 5\,V.
  \item \textbf{ASM-064} (Radio TM-D710GA/TS-2000): [AUDIO/DATA] TNC$\leftrightarrow$Radio, [RF] Antena, [PWR] 12\,V.
  \item \textbf{ASM-065} (Antenas UHF): [RF] Radio$\leftrightarrow$Antena, [MEC] montaje en STR.
\end{itemize}
(\emph{Ref. extensión PBS/ASM/N2 del SyRS})%

\subsubsection{Verificación de la interfaz (ICD checks)}
\begin{enumerate}
  \item \textbf{UDP/IP + CSP (Wi-Fi)}: captura pcap en GCS/rover con decodificación CSP; latencia inyectada 0–5\,s (CON-002).
  \item \textbf{KISS/AX.25 + CSP (UHF)}: consola serie KISS + logs del TNC; verificación de CRC-16 AX.25 y decodificación CSP.
  \item \textbf{Rutas CSP}: prueba directa y vía Relay (store-and-forward) con evidencias de recepción/reenvío/ACK.
  \item \textbf{RDP}: demostración de entrega $\geq 95\%$ con parámetros de campaña y registros de ACK/retry.
\end{enumerate}
(\emph{Ref. V\&V-COMMS-01/02 del SyRS})%

\subsubsection{Parámetros a cerrar (TBD) y valores resueltos}
\begin{itemize}
  \item \textbf{Puertos UDP} CMD/TM: \textbf{9000/9001} (\emph{resuelto — \texttt{config.py}}).
  \item \textbf{Puertos CSP} de servicio CMD/TM/HB: \textbf{11/10/1} (\emph{resuelto — \texttt{config.py}}).
  \item \textbf{IDs CSP de nodo} GCS/HLC: \textbf{1/2} (\emph{resuelto — \texttt{config.py}}).
  \item \textbf{Direcciones IP} (GCS/ROVER): \textbf{TBD} (asignar en run card de campaña).
  \item \textbf{Baudios KISS} (OBC$\leftrightarrow$TNC): \textbf{TBD} (\emph{rec.: 115200}) — aplica solo en Nivel~3.
  \item \textbf{timeout\_RDP / max\_retries}: \textbf{TBD} (\emph{rec.: 500\,ms / 3}) — aplica solo en Nivel~2+.
  \item \textbf{Cableado RF/PTT/audio}: \textbf{TBD} según radio/antena definitivos — aplica solo en Nivel~3.
\end{itemize}

\subsubsection{Notas}
El canal Wi-Fi se destina a pruebas y operación cercana; el canal UHF prioriza TM/CMD de misión cuando
el relieve penaliza la LoS. La selección dual-homing (UDP vs KISS) se decide por disponibilidad de enlace
y política RDP.%

\subsubsection{Cumplimiento con SyRS}
Este ICD satisface las secciones §9.1.1/\allowbreak §9.4.1 y los requisitos RF-006, COMM-REQ-001/\allowbreak 002/\allowbreak 003/\allowbreak 004/\allowbreak 005;
vincula N2/\allowbreak PBS/\allowbreak ASM y los casos de V\&V (V\&V-COMMS-01/\allowbreak 02).%
